\section{Related Work}

\para{Downstream applications of \molm and \smi.}
Both encoders we study are widely used as fixed feature extractors or finetuning backbones for downstream chemistry tasks. \molm has been applied to molecular property prediction across the MoleculeNet benchmarks~\citep{ross2022largescalechemicallanguagerepresentations,wu2018moleculenet} and adapted to ADMET and bioactivity prediction with parameter-efficient finetuning~\citep{harod2025effichem}. \smi has been used as a backbone for multi-scale battery materials and electrolyte discovery~\citep{sharma2025foundation}, and SMILES encoders of this family appear as components in multimodal property predictors that combine text and graph features~\citep{rollins2024molprop} and in scaling studies of chemical foundation models~\citep{cai2026chemfm}. These applications treat the encoders as black-box feature sources, which motivates our question of where in their layer stacks chemical information is actually represented and used.

\para{Interpreting scientific foundation models.}
On chemistry, \citet{pinto2025superior} find that intermediate layers of SMILES encoders outperform the final layer for downstream tasks, \citet{cohen2025unveiling} apply sparse autoencoders to \smi and steer generation by intervening on recovered features, and \citet{varadi2025circuits,muller2025uncovering,hu2024empirical} interpret attention heads, prediction circuits, and fragment-level features of other SMILES models. A parallel line of work targets biological foundation models: \citet{lu2026mechanisms} use counterfactual interventions on ESMFold latents to trace how the model builds spatial structure from sequence, \citet{goodfire2026evee} probe Evo~2 to generate interpretable predictions for millions of genetic variants, and \citet{goodfire2026alzheimers} use sparse autoencoders on an epigenetic foundation model to identify a novel Alzheimer's biomarker class. Building on this line of work, we pair layerwise linear probing with causal intervention on probe directions in \molm and \smi, which lets us separate information that is merely linearly recoverable from information the encoder actually uses during computation.
